\chapter{Réussir l'épreuve de mathématiques}

\noindent L'épreuve de mathématiques du Baccalauréat~C dure \textbf{4~heures} et compte
parmi les plus fortes : \textbf{coefficient~7}. Connaître son cours ne suffit pas — il
faut aussi \emph{savoir composer}. Ce chapitre te donne la méthode.

% ════════════════════════════════════════════════════════════════
\section{Comprendre l'épreuve (approche par compétences)}

L'épreuve suit l'\textbf{APC} (approche par compétences) et se compose de deux parties.

\begin{aretenir}
\textbf{Partie A — Évaluation des ressources} ($\approx$13 à 15~pts) : plusieurs
exercices indépendants qui testent tes \emph{savoirs} et \emph{savoir-faire}
(arithmétique, analyse, géométrie, probabilités…).\\[4pt]
\textbf{Partie B — Évaluation des compétences} ($\approx$5 à 7~pts) : \emph{une}
situation-problème de la vie courante, où tu dois \emph{mobiliser} tes connaissances
pour répondre à des tâches. \textbf{0,5~point} est réservé à la \emph{présentation}.
\end{aretenir}

\begin{remarque}
La répartition exacte des points varie d'une session à l'autre (par exemple 13,25/6,75
en 2023). Lis toujours le barème en tête de sujet : il t'indique où sont les points.
\end{remarque}

\section{Gérer les 4 heures}

\begin{enumerate}[label=\textbf{\arabic*.}]
\item \textbf{0–10 min : lis tout le sujet.} Repère les questions que tu sais faire
      tout de suite et celles qui demandent réflexion. Numérote-les mentalement par ordre d'attaque.
\item \textbf{Attaque par le plus sûr.} Commence par les exercices où tu es à l'aise :
      tu sécurises des points et tu te mets en confiance. L'ordre du sujet n'est pas obligatoire.
\item \textbf{Répartis le temps.} Compte environ \textbf{12 min par point} : un exercice
      sur 4~points mérite $\approx$45~min, pas plus. Si tu bloques, passe et reviens.
\item \textbf{Garde 40 min pour la Partie B.} La situation-problème rapporte gros et
      beaucoup d'élèves la bâclent faute de temps : ne tombe pas dans ce piège.
\item \textbf{Réserve 10 min de relecture.} Vérifie les calculs, les signes, les unités,
      et que tu as \emph{conclu} chaque question.
\end{enumerate}

\section{Bien rédiger une copie}

\begin{itemize}
\item \textbf{Justifie tout.} En maths, une réponse sans justification ne vaut presque
      rien. Cite le théorème ou la propriété utilisée, et \emph{vérifie ses hypothèses}.
\item \textbf{Structure.} Numérote tes réponses comme le sujet (1.a, 1.b, …). Saute des
      lignes entre les questions. Une copie aérée se corrige mieux — et se note mieux.
\item \textbf{Encadre tes résultats.} Le correcteur doit trouver ta conclusion d'un coup d'œil.
\item \textbf{Rédige des phrases.} « Donc $f$ est croissante sur $[0,+\infty[$ » vaut mieux
      qu'une flèche isolée. Introduis tes calculs (« Calculons… », « On en déduit… »).
\item \textbf{Soigne les figures.} Grandes, à la règle, légendées, avec le repère et les
      points nommés. Une bonne figure rapporte des points et évite des erreurs.
\end{itemize}

\section{Réussir la situation-problème (Partie B)}

\begin{aretenir}
Méthode en 4 temps pour chaque tâche : \textbf{(1) Comprendre} ce qui est demandé ;
\textbf{(2) Modéliser} (poser les inconnues, traduire l'énoncé en équations / fonction /
suite / probabilité) ; \textbf{(3) Résoudre} avec les outils du cours ; \textbf{(4) Conclure}
par une \emph{phrase} qui répond à la question posée, avec l'unité.
\end{aretenir}

\noindent Le correcteur évalue ta \emph{démarche} autant que le résultat : même si tu
n'aboutis pas, une modélisation correcte et des étapes justifiées rapportent des points.
N'oublie pas les \textbf{0,5~point de présentation} : copie propre, soignée, organisée.

\section{Les erreurs à ne pas commettre}

\subsection{Erreurs générales}
\begin{itemize}
\item Oublier de \textbf{justifier} (théorème non cité, hypothèses non vérifiées).
\item Confondre \textbf{condition nécessaire} et \textbf{suffisante} ; « $\Rightarrow$ » et « $\iff$ ».
\item Ne pas \textbf{conclure} : un calcul juste sans phrase de conclusion perd des points.
\item Mauvaise \textbf{gestion du temps} : rester bloqué 1~h sur une question.
\end{itemize}

\subsection{Analyse}
\begin{itemize}
\item Oublier le \textbf{domaine de définition} (ex. $\ln$ exige l'argument $>0$ ; $\sqrt{\ }$ exige $\ge0$).
\item \textbf{Dériver un quotient} en oubliant le carré au dénominateur ou en inversant : $\left(\frac uv\right)'=\dfrac{u'v-uv'}{v^2}$.
\item Conclure les variations sans \textbf{étudier le signe de $f'$} (un tableau de signes propre).
\item Appliquer le \textbf{T.V.I.} pour l'\emph{unicité} sans vérifier la \textbf{stricte monotonie}.
\item Erreur de \textbf{croissances comparées} : en $+\infty$, $\mathrm{e}^x$ « l'emporte » sur toute puissance, qui l'emporte sur $\ln$.
\item Oublier la \textbf{constante} d'intégration, ou les bornes lors d'un changement.
\end{itemize}

\subsection{Suites}
\begin{itemize}
\item Confondre suite \textbf{arithmétique} ($+r$) et \textbf{géométrique} ($\times q$).
\item Affirmer qu'une suite converge sans justifier (\emph{croissante majorée}, etc.).
\item Conclure $u_n\to\ell$ \emph{donc} $\ell=f(\ell)$ sans dire que $f$ est continue.
\item Erreur d'indice dans une somme : $\sum_{k=0}^{n}q^k=\dfrac{1-q^{n+1}}{1-q}$ (et non $q^n$).
\end{itemize}

\subsection{Arithmétique}
\begin{itemize}
\item Appliquer le \textbf{théorème de Gauss} sans l'hypothèse $a\wedge b=1$.
\item Oublier de vérifier que $\mathrm{pgcd}(a,b)\mid c$ avant de résoudre $ax+by=c$.
\item Donner \emph{une} solution d'une équation diophantienne au lieu de \textbf{toutes}.
\item Erreur de \textbf{congruence} : $a\equiv b\ [n]$ se manipule comme une égalité pour $+,\times$, mais pas pour la division.
\end{itemize}

\subsection{Nombres complexes \& géométrie}
\begin{itemize}
\item Oublier qu'un argument est défini \textbf{modulo $2\pi$}.
\item Confondre \textbf{module} et \textbf{argument} ; mal placer un point dans le plan.
\item Pour une similitude $z'=az+b$ : oublier que le centre n'existe que si $a\ne1$, $\omega=\frac{b}{1-a}$.
\item Barycentre : oublier la condition $\sum\alpha_i\ne0$.
\item Produit scalaire : confondre $\vec u\cdot\vec v$ (scalaire) et $\vec u\wedge\vec v$ (vecteur).
\end{itemize}

\subsection{Probabilités}
\begin{itemize}
\item Confondre tirage \textbf{avec} / \textbf{sans} remise (et donc loi binomiale / hypergéométrique).
\item Oublier que les probabilités d'une loi doivent \textbf{sommer à 1} (sers-t'en pour vérifier !).
\item « Au moins un » : penser à l'\textbf{événement contraire} $1-P(\text{aucun})$.
\item Mal poser une probabilité conditionnelle : $P_B(A)=\dfrac{P(A\cap B)}{P(B)}$.
\end{itemize}

\begin{synthese}[En une phrase]
Lis tout, attaque le plus sûr, justifie chaque étape, soigne figures et présentation,
garde du temps pour la Partie~B, et \emph{conclus} toujours. C'est ainsi qu'on transforme
des connaissances en \textbf{points}.
\end{synthese}
